% ============================================================
%  ch2_background_concepts.tex
%  Chapter 2: Background and Concepts
%  Master's Thesis -- SQL MATCH_RECOGNIZE on Pandas DataFrames
% ============================================================
% !TEX root = Thesis.tex
\chapter{Background and Concepts}
\label{chap:background}

This chapter introduces the foundations needed to understand the
\texttt{MATCH\_RECOGNIZE} engine for Pandas DataFrames.
Section~\ref{sec:bg-pattern-matching} explains pattern matching in
sequential data.  Section~\ref{sec:bg-match-recognize} introduces the
SQL:2016 clause and the terms used in this thesis.  A V-pattern example
then connects these concepts to the proposed system.  The next chapter
reviews related work and identifies the research gap.

% ----------------------------------------------------------------
\section{Pattern Matching in Sequential Data Analysis}
\label{sec:bg-pattern-matching}
% ----------------------------------------------------------------

Sequential pattern detection identifies ordered events or states in a
dataset.  A conventional filter tests each row against one or more
predicates.  A row-pattern query also tests the order and relationships
between rows.  It returns a group when the rows follow the requested
sequence.  This makes it possible to find shapes such as V, W, and U,
as well as repeated or changing behavior.  Common applications include
security monitoring, stock-price analysis, sensor analysis, fraud
detection, and complex event processing~\cite{OracleMR}.

The important distinction is that the ``symbols'' in row-pattern
recognition are not characters in a string.  They are data rows that
become associated with pattern variables such as \texttt{A},
\texttt{B}, or \texttt{DOWN} when their predicates hold.  A row
sequence can therefore be read as a sequence of variable assignments.
This explains why regular-expression-like syntax and finite automata
fit the problem.  Later chapters use this connection to compile row
patterns into automata and execute them over ordered Pandas partitions.

Relational databases are strong at set-based operations such as
selection, projection, joins, and aggregation.  Sequence analysis was
traditionally harder.  Analysts often exported data to a stream
processor, used cursors and temporary tables, or wrote the logic in an
application.  SQL:2016 introduced row-pattern recognition so that this
logic could be expressed directly in SQL~\cite{ISO2016SQL}.

% ................................
\subsection{Real-World Applications}
\label{subsec:bg-applications}
% ................................

\paragraph{Financial Market Analysis.}
Detecting a V-shaped price pattern---a decline followed by a
recovery---requires an ordered sequence of price movements.  Successive
higher troughs, for example, can indicate an upward trend.  The worked
example later in this chapter applies a V-pattern to ordered price
records~\cite{OracleMR}.

\paragraph{Fraud Detection.}
Credit card fraud can appear as a sequence, such as several small test
charges before a large purchase.  Another example is a series of
transactions in locations that are too far apart to be realistic.
Writing these rules in SQL can make them easier to manage than
procedural code~\cite{petkovic2022fraud}.

\paragraph{IoT and Sensor Data.}
Manufacturing systems monitor sensor readings for signs of equipment
degradation.  One example is a gradual temperature rise followed by a
sudden drop.  Pattern queries over ordered sensor readings are a
natural use of \texttt{MATCH\_RECOGNIZE}~\cite{OracleMR}.

\paragraph{Clickstream and Funnel Analysis.}
E-commerce platforms study the paths that lead to a purchase or an
exit.  A funnel query may look for
``view~$\rightarrow$ add-to-cart~$\rightarrow$ remove~$\rightarrow$
exit.''  With \texttt{MATCH\_RECOGNIZE}, one query can count repeated
actions, return their timestamps, and measure the elapsed time.  The
same analysis otherwise needs several subqueries or custom session
logic~\cite{hoffa2021snowflake}.

\paragraph{Network Security and Intrusion Detection.}
Cyber-attacks often appear as a series of events rather than one
isolated event.  A multi-stage intrusion may start with port scanning,
continue with failed logins, and end with a successful login and
privilege escalation.  One row-pattern query can describe this
sequence and return the relevant session, timestamps, port, and
account~\cite{petkovic2022fraud}.

\paragraph{Healthcare and Patient Monitoring.}
Early-warning systems can look for a sequence of changes in clinical
measurements.  An alert may combine a rising heart rate, falling blood
pressure, and a later laboratory change.  A row-pattern query can
describe this ordered sequence.  The actual alerting rule depends on
the clinical domain and system.

\paragraph{Supply Chain and Logistics.}
One late shipment may lead to partial deliveries and then a stockout.
A row-pattern query over a fulfillment log can find this sequence for
each supplier or warehouse.  It can also measure the time from the
first delay to the final stockout.

\paragraph{Financial Compliance and Trade Surveillance.}
Trade-surveillance systems look for sequences that may indicate market
manipulation.  In a layering example, a trader places a large resting
order, trades on the other side, and then cancels the first order.
\texttt{MATCH\_RECOGNIZE} can express this timed sequence as an
auditable SQL query.

% ----------------------------------------------------------------
\section{The SQL:2016 \texttt{MATCH\_RECOGNIZE} Clause}
\label{sec:bg-match-recognize}
% ----------------------------------------------------------------

The SQL:2016 standard introduced
\texttt{MATCH\_RECOGNIZE}~\cite{ISO2016SQL,iso2016}.  It finds ordered
rows in which each part of a pattern is defined by a row condition.
The regular-expression-like pattern describes event sequences in a
clear, declarative form~\cite{winand2017sqlnew}.  Similar logic can be
written with joins, window functions, recursion, or procedural code.
The clause makes that logic shorter and easier to maintain.

The SQL standard distinguishes three optional features for row pattern
recognition~\cite{winand_match_recognize}.  \textbf{Feature~R010}
covers row pattern recognition in the \texttt{FROM} clause.
\textbf{Feature~R020} covers row patterns in \texttt{WINDOW}
specifications.  \textbf{Feature~R030} covers full aggregate support.
They are distinct parts of the standard, not successive
levels~\cite{ISO2021,michels2018new,melton2017sqlrpr}.

This thesis uses the SQL:2016 terminology as the conceptual frame,
but the implementation claims are intentionally scoped.  The proposed
engine targets core R010-style behavior in the \texttt{FROM} clause
and evaluates selected row-pattern constructs that are relevant to
in-memory Pandas analytics.  Features outside this scope, including
full window-clause row-pattern recognition and complete vendor-specific
extensions, are treated as related work or future work rather than as
requirements of the implemented system.


% ................................
\subsection{Clause Structure and Components}
\label{subsec:bg-clause-structure}
% ................................
The clause works with a table or a derived table.  It is commonly
described through the eight specifications shown in
Listing~\ref{lst:mr-syntax}.  \texttt{PATTERN} is required.  In the
evaluated engine, \texttt{DEFINE} and the other specifications are
optional.  Together, they control partitioning, order, output, row
classification, and overlap between matches.

\begin{lstlisting}[
    language     = SQL,
    caption      = {General clause structure used in this thesis},
    label        = {lst:mr-syntax}
]
SELECT <columns>
FROM   <table>
MATCH_RECOGNIZE (
    [PARTITION BY <partition_columns>] -- 1. Data organization
    [ORDER BY     <ordering_columns>]  -- 2. Row sequencing
    [MEASURES     <expressions>]       -- 3. Output computations
    [<rows_per_match>]                 -- 4. Result granularity
    [AFTER MATCH SKIP <option>]        -- 5. Overlap handling
    PATTERN ( <pattern_regex> )        -- 6. Sequence specification
    [SUBSET <union_variable_defs>]     -- 7. Union variables
    [DEFINE <variable_definitions>]    -- 8. Variable predicates
) [AS <alias>]
\end{lstlisting}

\begin{enumerate}

\item \textbf{PARTITION BY.}
If present, this specification splits the input into independent
partitions, similar to a window function.  Matching then runs within
each customer, sensor, or other partition.  Independent partitions may
also support parallel processing.  If it is omitted, the input forms
one partition.

\item \textbf{ORDER BY.}
This determines the row order within each partition.  The matcher
processes rows in that sequence.  An explicit and deterministic order
is normally needed for repeatable results.  Several sort keys may be
combined, each with its own \texttt{ASC} or \texttt{DESC} direction.
This describes the standard syntax.  The evaluated R010 execution path
in this thesis is narrower and uses ascending pattern order.


\item \textbf{MEASURES.}
  Optionally defines computed expressions that extract values from
  matched row sequences using pattern variables and navigation functions
  (\texttt{FIRST}, \texttt{LAST}, \texttt{PREV}, \texttt{NEXT}).
  Each measure expression becomes a column in the output, allowing a
  match to be summarized without returning every raw input row.

\item \textbf{Output options.}
  \begin{itemize}
    \item \texttt{ONE ROW PER MATCH}: one summary row per complete
      match.
    \item \texttt{ALL ROWS PER MATCH}: one output row for each included
      row in a match.
    \item \texttt{SHOW EMPTY MATCHES}: include an output row when a
      pattern completes without consuming an input row.
    \item \texttt{OMIT EMPTY MATCHES}: suppress such empty-match rows.
    \item \texttt{WITH UNMATCHED ROWS}: an
      \texttt{ALL ROWS PER MATCH} modifier that includes
      non-matching rows in the result where supported.
  \end{itemize}

\item \textbf{AFTER MATCH SKIP.}
  Controls how the engine resumes after finding a match:
  \begin{itemize}
    \item \texttt{SKIP TO NEXT ROW} -- resume from the row after
      match start (maximum overlaps).
    \item \texttt{SKIP PAST LAST ROW} -- resume from the row after
      match end (non-overlapping matches).
    \item \texttt{SKIP TO FIRST} \textit{var} -- resume from the
      first row classified as \textit{var}.
    \item \texttt{SKIP TO LAST} \textit{var} -- resume from the last
      row classified as \textit{var}.
  \end{itemize}

\item \textbf{PATTERN.}
  Specifies the sequence using regular-expression-like syntax with
  pattern variables as terminals:
  \begin{itemize}
    \item Concatenation: \texttt{A~B~C} (A followed by B followed by C).
    \item Alternation: \texttt{A~|~B} (either A or B).
    \item Quantifiers: \texttt{A*} (zero or more), \texttt{A+}
      (one or more), \texttt{A?} (zero or one), \texttt{A\{n\}}
      (exactly $n$), \texttt{A\{n,m\}} (between $n$ and $m$).
      A trailing question mark gives the reluctant form, such as
      \texttt{A*?}, \texttt{A+?}, \texttt{A??}, or
      \texttt{A\{n,m\}?}.
    \item Grouping: \texttt{(A~B)+} (one or more AB pairs).
    \item Anchors: \texttt{\^{}A~B~C\$} (pattern must span the
      entire partition).
    \item Exclusions: \texttt{\{-~A+~B+~-\}} (rows matching the
      enclosed sub-pattern are excluded from row-level output).
    \item Permutations: \texttt{PERMUTE(A,~B,~C)} (any ordering
      of A, B, and C).
  \end{itemize}

\item \textbf{SUBSET.}
This optional specification defines a union variable from variables
already used in \texttt{PATTERN}.  It does not change the match.
Instead, it gives one name to rows classified by any of its component
variables.  It can then be referenced in \texttt{MEASURES},
\texttt{DEFINE}, and supported skip targets.


\item \textbf{DEFINE.}
This optional specification assigns Boolean predicates to pattern
variables.  The matcher tests a candidate row in the current match
context to decide whether it can be classified as a variable.
Predicates may use current-row columns, navigation functions,
aggregations over matched rows, and Boolean or comparison
operators~\cite{winand_match_recognize}.  A variable without a
\texttt{DEFINE} predicate uses the implicit true predicate.

\end{enumerate}

The clause forms a table expression inside a \texttt{SELECT} query.
An outer query can therefore filter, group, aggregate, or order its
result where the SQL dialect allows~\cite{winand_match_recognize}.

These specifications separate the sequence to detect from the way
matches are classified and returned.  \texttt{PATTERN} defines the
sequence, while \texttt{DEFINE} gives the row-classification rules.
The output and skip options decide how matches are shown and where the
next search begins.  This separation guides the system design in later
chapters: parsing creates a structured query, validation checks its
variables and columns, and compilation builds an automaton for
execution.


% ................................
\subsection{Illustrative Example: Price V-Pattern}
\label{subsec:bg-example}
% ................................

Listing~\ref{lst:mr-example} shows a \texttt{MATCH\_RECOGNIZE}
query that detects a V-shaped sequence in customer order prices.

\begin{lstlisting}[
    language = SQL,
    caption  = {Detecting V-shaped price patterns with \texttt{MATCH\_RECOGNIZE}},
    label    = {lst:mr-example}
]
SELECT customer_id,
       start_price, bottom_price, final_price,
       start_date,  final_date
FROM   orders
MATCH_RECOGNIZE (
    PARTITION BY customer_id
    ORDER BY     order_date
    MEASURES
        START.price          AS start_price,
        LAST(DOWN.price)     AS bottom_price,
        LAST(UP.price)       AS final_price,
        START.order_date     AS start_date,
        LAST(UP.order_date)  AS final_date
    ONE ROW PER MATCH
    AFTER MATCH SKIP PAST LAST ROW
    PATTERN (START DOWN+ UP+)
    DEFINE
        DOWN AS price < PREV(price),
        UP   AS price > PREV(price)
);
\end{lstlisting}


The query groups rows by \texttt{customer\_id} and orders each group by
\texttt{order\_date}.  A match starts with any row.  It then has one or
more price decreases followed by one or more price increases.  For each
match, \texttt{MEASURES} returns the starting value and date, the value
at the last \texttt{DOWN} row, and the final \texttt{UP} value and date.

% ----------------------------------------------------------------
\section{Chapter Summary}
\label{sec:bg-summary}
% ----------------------------------------------------------------
This chapter explained sequential pattern matching and why ordinary
row-by-row predicates are difficult to use for ordered events.  It
covered the main \texttt{MATCH\_RECOGNIZE} specifications, including
partitioning, ordering, measures, output modes, skip modes, patterns,
subsets, and variable predicates.

The V-pattern example showed how these parts fit together in one query
and why the abstraction works with automata-based execution.
Chapter~\ref{chap:related} next compares existing implementations,
their deployment models, and their practical limits.
